% GENERATED FILE. Edit canonical lessons, then run npm run generate:books.
% canonical-source-hash: fnv1a64:de2014178d196649
% canonical-lessons: HI-C389-behtar, HI-C389-badtar, HI-C389-barhiya, HI-C389-samanya, HI-C389-nikharna

\chapter{Better and Worse}
\label{ch:hi-better-and-worse}
\hlchaptermodality{389}

\begin{chapteropening}
\textbf{By the end of this chapter:} \emph{I can say whether things will get better or worse.}

Complete the last lesson of chapter 389: I can say whether things will get better or worse.
\end{chapteropening}

\section[\texorpdfstring{\hi{बेहतर} (behtar)}{बेहतर (behtar)}]{\hi{बेहतर} (behtar) — better}

\label{lesson:HI-C389-behtar}



\begin{quote}
Before the new one: say the Hindi for minor, then the Hindi for close to.
\end{quote}

\subsection*{You'll want to know: \hi{बेहतर}}
\textbf{\hi{बेहतर}} — \emph{behtar} — \textquotedblleft{}better\textquotedblright{}.

\begin{quote}
\hi{अगर} \hi{तुम} \hi{आराम} \hi{करोगे} \hi{तो} \hi{कल} \hi{तबीयत} \hi{बेहतर} \hi{होगी।}
\end{quote}

\begin{quote}
\emph{agar tum ārām karoge to kal tabīyat behtar hogī.}
\end{quote}

\begin{quote}
If you rest, you'll feel better tomorrow.
\end{quote}

\begin{cousinweb}
\textbf{\hi{बेहतर}} (\emph{behtar}) is usually traced to Persian \emph{behtar}, \textquotedblleft{}better\textquotedblright{}.
\end{cousinweb}

\begin{grammarlens}[title={What you've built}]
One more word for saying what happens if.
\end{grammarlens}

\subsection*{Your turn}
\begin{itemize}
\raggedright
  \item \emph{Say it:} \emph{behtar}
  \item \emph{Say it:} \emph{behtar}, once more
  \item \emph{Say it:} \emph{agar tum ārām karoge to kal tabīyat behtar hogī.} — If you rest, you'll feel better tomorrow.
  \item \emph{Recall:} say the Hindi for minor, then the Hindi for close to, then say \emph{behtar} again
\end{itemize}

\begin{tcolorbox}[breakable,colback=teal!4,colframe=teal!35!black,title={Before you move on}]
What does \hi{बेहतर} mean? (\textquotedblleft{}Better\textquotedblright{}.) Say it once more.
\end{tcolorbox}
\section[\texorpdfstring{\hi{बदतर} (badtar)}{बदतर (badtar)}]{\hi{बदतर} (badtar) — worse}

\label{lesson:HI-C389-badtar}



\begin{quote}
Before the new one: say the Hindi for close to, then the Hindi for better.
\end{quote}

\subsection*{You'll want to know: \hi{बदतर}}
\textbf{\hi{बदतर}} — \emph{badtar} — \textquotedblleft{}worse\textquotedblright{}.

\begin{quote}
\hi{अगर} \hi{इलाज} \hi{नहीं} \hi{हुआ} \hi{तो} \hi{हालत} \hi{बदतर} \hi{हो} \hi{जाएगी।}
\end{quote}

\begin{quote}
\emph{agar ilāj nahīṃ huā to hālat badtar ho jāegī.}
\end{quote}

\begin{quote}
If there's no treatment, the condition will get worse.
\end{quote}

\begin{cousinweb}
\textbf{\hi{बदतर}} (\emph{badtar}) is usually traced to Persian \emph{badtar}, \textquotedblleft{}worse\textquotedblright{}.
\end{cousinweb}

\begin{grammarlens}[title={What you've built}]
One more word for saying what happens if.
\end{grammarlens}

\subsection*{Your turn}
\begin{itemize}
\raggedright
  \item \emph{Say it:} \emph{badtar}
  \item \emph{Say it:} \emph{badtar}, once more
  \item \emph{Say it:} \emph{agar ilāj nahīṃ huā to hālat badtar ho jāegī.} — If there's no treatment, the condition will get worse.
  \item \emph{Recall:} say the Hindi for close to, then the Hindi for better, then say \emph{badtar} again
\end{itemize}

\begin{tcolorbox}[breakable,colback=teal!4,colframe=teal!35!black,title={Before you move on}]
What does \hi{बदतर} mean? (\textquotedblleft{}Worse\textquotedblright{}.) Say it once more.
\end{tcolorbox}
\section[\texorpdfstring{\hi{बढ़िया} (baṛhiyā)}{बढ़िया (baṛhiyā)}]{\hi{बढ़िया} (baṛhiyā) — excellent, first-rate}

\label{lesson:HI-C389-barhiya}



\begin{quote}
Before the new one: say the Hindi for better, then the Hindi for worse.
\end{quote}

\subsection*{You'll want to know: \hi{बढ़िया}}
\textbf{\hi{बढ़िया}} — \emph{baṛhiyā} — \textquotedblleft{}excellent, first-rate\textquotedblright{}.

\begin{quote}
\hi{अगर} \hi{तुम} \hi{यह} \hi{पराठा} \hi{खाओगे} \hi{तो} \hi{कहोगे} \hi{कि} \hi{बहुत} \hi{बढ़िया} \hi{है।}
\end{quote}

\begin{quote}
\emph{agar tum yah parāṭhā khāoge to kahoge ki bahut baṛhiyā hai.}
\end{quote}

\begin{quote}
If you eat this paratha, you'll say it's excellent.
\end{quote}

\begin{cousinweb}
\textbf{\hi{बढ़िया}} (\emph{baṛhiyā}) is a tadbhava, inherited from Sanskrit \emph{vardhita}, \textquotedblleft{}grown, increased\textquotedblright{} (CDIAL 11384).
\end{cousinweb}

\begin{grammarlens}[title={What you've built}]
One more word for saying what happens if.
\end{grammarlens}

\subsection*{Your turn}
\begin{itemize}
\raggedright
  \item \emph{Say it:} \emph{baṛhiyā}
  \item \emph{Say it:} \emph{baṛhiyā}, once more
  \item \emph{Say it:} \emph{agar tum yah parāṭhā khāoge to kahoge ki bahut baṛhiyā hai.} — If you eat this paratha, you'll say it's excellent.
  \item \emph{Recall:} say the Hindi for better, then the Hindi for worse, then say \emph{baṛhiyā} again
\end{itemize}

\begin{tcolorbox}[breakable,colback=teal!4,colframe=teal!35!black,title={Before you move on}]
What does \hi{बढ़िया} mean? (\textquotedblleft{}Excellent, first-rate\textquotedblright{}.) Say it once more.
\end{tcolorbox}
\section[\texorpdfstring{\hi{सामान्य} (sāmānya)}{सामान्य (sāmānya)}]{\hi{सामान्य} (sāmānya) — normal, usual}

\label{lesson:HI-C389-samanya}



\begin{quote}
Before the new one: say the Hindi for worse, then the Hindi for first-rate.
\end{quote}

\subsection*{You'll want to know: \hi{सामान्य}}
\textbf{\hi{सामान्य}} — \emph{sāmānya} — \textquotedblleft{}normal, usual\textquotedblright{}.

\begin{quote}
\hi{अगर} \hi{बुख़ार} \hi{सामान्य} \hi{से} \hi{ज़्यादा} \hi{होगा} \hi{तो} \hi{डाक्टर} \hi{को} \hi{बुलाओ।}
\end{quote}

\begin{quote}
\emph{agar bukhār sāmānya se zyādā hogā to ḍākṭar ko bulāo.}
\end{quote}

\begin{quote}
If the fever is above normal, call the doctor.
\end{quote}

\begin{cousinweb}
\textbf{\hi{सामान्य}} (\emph{sāmānya}) is a tatsama, borrowed from Sanskrit \emph{sāmānya}, \textquotedblleft{}common, general\textquotedblright{} (Monier-Williams).
\end{cousinweb}

\begin{grammarlens}[title={What you've built}]
One more word for saying what happens if.
\end{grammarlens}

\subsection*{Your turn}
\begin{itemize}
\raggedright
  \item \emph{Say it:} \emph{sāmānya}
  \item \emph{Say it:} \emph{sāmānya}, once more
  \item \emph{Say it:} \emph{agar bukhār sāmānya se zyādā hogā to ḍākṭar ko bulāo.} — If the fever is above normal, call the doctor.
  \item \emph{Recall:} say the Hindi for worse, then the Hindi for first-rate, then say \emph{sāmānya} again
\end{itemize}

\begin{tcolorbox}[breakable,colback=teal!4,colframe=teal!35!black,title={Before you move on}]
What does \hi{सामान्य} mean? (\textquotedblleft{}Normal, usual\textquotedblright{}.) Say it once more.
\end{tcolorbox}
\section[\texorpdfstring{\hi{निखरना} (nikharnā)}{निखरना (nikharnā)}]{\hi{निखरना} (nikharnā) — to brighten, to come out at its best}

\label{lesson:HI-C389-nikharna}



\begin{quote}
Before the new one: say the Hindi for first-rate, then the Hindi for normal.
\end{quote}

\subsection*{You'll want to know: \hi{निखरना}}
\textbf{\hi{निखरना}} — \emph{nikharnā} — \textquotedblleft{}to brighten, to come out at its best\textquotedblright{}.

\begin{quote}
\hi{अगर} \hi{तुम} \hi{रोज़} \hi{अभ्यास} \hi{करोगे} \hi{तो} \hi{तुम्हारा} \hi{गाना} \hi{निखर} \hi{जाएगा।}
\end{quote}

\begin{quote}
\emph{agar tum roz abhyās karoge to tumhārā gānā nikhar jāegā.}
\end{quote}

\begin{quote}
If you practise every day, your singing will come into its own.
\end{quote}

Past (it …): \hi{वह} \hi{रंग} \hi{धूप} \hi{में} \hi{निखर} \hi{गया}, \emph{vah raṅg dhūp meṁ nikhar gayā}.

\begin{cousinweb}
\textbf{\hi{निखरना}} (\emph{nikharnā}) is inherited from \emph{*niḥkṣarati}, \textquotedblleft{}flows out\textquotedblright{}, a form CDIAL reconstructs from its descendants (CDIAL 7095). CDIAL glosses the Hindi verb \textquotedblleft{}to become clear\textquotedblright{}.
\end{cousinweb}

\begin{grammarlens}[title={What you've built}]
One more word for saying what happens if.
\end{grammarlens}

\subsection*{Your turn}
\begin{itemize}
\raggedright
  \item \emph{Say it:} \emph{nikharnā}
  \item \emph{Say it:} \emph{nikharnā}, once more
  \item \emph{Say it:} \emph{agar tum roz abhyās karoge to tumhārā gānā nikhar jāegā.} — If you practise every day, your singing will come into its own.
  \item \emph{Recall:} say the Hindi for first-rate, then the Hindi for normal, then say \emph{nikharnā} again
\end{itemize}

\begin{tcolorbox}[breakable,colback=teal!4,colframe=teal!35!black,title={Before you move on}]
What does \hi{निखरना} mean? (\textquotedblleft{}To brighten, to come out at its best\textquotedblright{}.) Say it once more.
\end{tcolorbox}
